\documentclass[aspectratio=169]{beamer}

\usetheme{Madrid}
\usecolortheme{seahorse}

\usepackage[utf8]{inputenc}
\usepackage{listings}
\usepackage{xcolor}
\usepackage{graphicx}

\definecolor{codegreen}{rgb}{0,0.6,0}
\definecolor{codegray}{rgb}{0.5,0.5,0.5}
\definecolor{codepurple}{rgb}{0.58,0,0.82}
\definecolor{backcolour}{rgb}{0.95,0.95,0.92}

\lstdefinestyle{cppstyle}{
    backgroundcolor=\color{backcolour},
    commentstyle=\color{codegreen},
    keywordstyle=\color{blue},
    numberstyle=\tiny\color{codegray},
    stringstyle=\color{codepurple},
    basicstyle=\ttfamily\small,
    breaklines=true,
    numbers=left,
    numbersep=5pt,
    tabsize=2,
    language=C++
}

\lstset{style=cppstyle}

\title{CS202 Seminar - Decorator Design Pattern}
\author{}
\date{}

\begin{document}

%%%%%%%%%%%%%%%%%%%%%%%%%%%%%%%%%%%%%%%%%%%%%%%%%%%%%

\begin{frame}
\titlepage
\end{frame}

%%%%%%%%%%%%%%%%%%%%%%%%%%%%%%%%%%%%%%%%%%%%%%%%%%%%%

\begin{frame}{Agenda}

\tableofcontents

\end{frame}

\section{Motivation}

%%%%%%%%%%%%%%%%%%%%%%%%%%%%%%%%%%%%%%%%%%%%%%%%%%%%%

\begin{frame}{1. A Real-world Problem}

Imagine your team is building an order management software for a coffee shop chain (e.g., The Coffee House).

Available coffee:

\begin{itemize}
\item Dark Roast
\item House Blend
\item Espresso
\end{itemize}

Customers rarely order plain coffee.

They often request additional condiments:

\begin{itemize}
\item Milk
\item Sugar
\item Caramel
\end{itemize}

Each topping changes

\begin{itemize}
\item Price
\item Description
\end{itemize}

\end{frame}

%%%%%%%%%%%%%%%%%%%%%%%%%%%%%%%%%%%%%%%%%%%%%%%%%%%%%

\section{Naive Solution}

\begin{frame}{2. Naive Solution - Inheritance}

The first idea is using inheritance.

\bigskip

Example:

\begin{itemize}

\item Beverage

\item DarkRoast

\item DarkRoastWithMilk

\item DarkRoastWithMilkAndSugar

\item DarkRoastWithMilkAndSugarAndWhip

\item EspressoWithWhip

\item HouseBlendWithMilk

\item ...

\end{itemize}

\pause

\begin{alertblock}{Problem}

Every combination requires a new class.

\end{alertblock}

\end{frame}

%%%%%%%%%%%%%%%%%%%%%%%%%%%%%%%%%%%%%%%%%%%%%%%%%%%%%

\begin{frame}[fragile]{Naive Solution - Boolean Flags}

Another approach is storing every topping inside the Beverage class.

\bigskip

\begin{lstlisting}
class Beverage {

private:

    bool hasMilk;
    bool hasSugar;
    bool hasWhip;
    bool hasCaramel;

public:
    ...
};
\end{lstlisting}

\pause

Everything is controlled by boolean variables.

\end{frame}

%%%%%%%%%%%%%%%%%%%%%%%%%%%%%%%%%%%%%%%%%%%%%%%%%%%%%

\section{Problems}

\begin{frame}{3. Problems with the Naive Solution}

\begin{block}{Class Explosion}

4 coffee types × many toppings = dozens (or hundreds) of subclasses.

\end{block}

\pause

\begin{block}{Hard Maintenance}

Changing the price of milk requires modifying many classes.

\end{block}

\pause

\begin{block}{Violates Open/Closed Principle}

Adding a new topping (e.g. Boba)

requires modifying Beverage.

\end{block}

\pause

\begin{block}{No Runtime Flexibility}

Cannot naturally support

\begin{itemize}
\item Double Milk
\item Triple Mocha
\item Dynamic combinations
\end{itemize}

\end{block}

\end{frame}

%%%%%%%%%%%%%%%%%%%%%%%%%%%%%%%%%%%%%%%%%%%%%%%%%%%%%

\section{Decorator Pattern}

\begin{frame}{4. Decorator Pattern}

Decorator is a \textbf{Structural Design Pattern}.

\bigskip

Instead of inheritance,

we dynamically attach new behaviors to an object.

\bigskip

Core philosophy:

\begin{center}

\Huge
Composition over Inheritance

\end{center}
\end{frame}

%%%%%%%%%%%%%%%%%%%%%%%%%%%%%%%%%%%%%%%%%%%%%%%%%%%%%

\begin{frame}{Core Concepts}

Decorator uses wrapper objects.

\bigskip

\begin{center}

Coffee

$\uparrow$

Milk Decorator

$\uparrow$

Sugar Decorator

$\uparrow$

Mocha Decorator

\end{center}
\end{frame}

%%%%%%%%%%%%%%%%%%%%%%%%%%%%%%%%%%%%%%%%%%%%%%%%%%%%%
\section{Class Diagram}

\begin{frame}{5. General Structure of Decorator Pattern}

\begin{block}{Main Components}

\begin{itemize}

\item \textbf{Component}: Common interface for all objects.

\item \textbf{ConcreteComponent}: The original object.

\item \textbf{Decorator}: Abstract wrapper containing a Component.

\item \textbf{ConcreteDecorator}: Adds new behavior while preserving the interface.

\end{itemize}

\end{block}

Flow:

\[
Component
\leftarrow
Decorator
\leftarrow
Decorator
\leftarrow
Decorator
\]

Every decorator is still recognized as a Component.

\end{frame}

%%%%%%%%%%%%%%%%%%%%%%%%%%%%%%%%%%%%%%%%%%%%%%%%%%%%%

\begin{frame}[fragile]{Coffee UML (Mermaid)}

\small

\begin{figure}
\centering
\includegraphics[width=0.4\textwidth]{diagram.jpg}
\caption{Coffee UML diagram}
\end{figure}

\end{frame}

%%%%%%%%%%%%%%%%%%%%%%%%%%%%%%%%%%%%%%%%%%%%%%%%%%%%%
\section{Implementation}

\begin{frame}[fragile]{7. Component}

\begin{lstlisting}
class Beverage {
public:
    virtual string getDescription() const = 0;
    virtual double cost() const = 0;
    virtual ~Beverage() = default;
};
\end{lstlisting}

\pause

\begin{block}{Role}

Defines the common interface.
Every coffee and every topping must inherit from Beverage.

\end{block}

\end{frame}

%%%%%%%%%%%%%%%%%%%%%%%%%%%%%%%%%%%%%%%%%%%%%%%%%%%%%

\begin{frame}[fragile]{Concrete Component}

\begin{lstlisting}
class DarkRoast : public Beverage {
public:
    string getDescription() const override {
        return "Dark Roast Coffee";
    }

    double cost() const override {
        return 30.0;
    }
};
\end{lstlisting}

\pause

DarkRoast is the original object before decoration.

\end{frame}

%%%%%%%%%%%%%%%%%%%%%%%%%%%%%%%%%%%%%%%%%%%%%%%%%%%%%

\begin{frame}[fragile]{Base Decorator}

\begin{lstlisting}
class CondimentDecorator : public Beverage {
protected:
    Beverage* beverage;
public:
    CondimentDecorator(Beverage* bev) : beverage(bev) {}

    virtual ~CondimentDecorator() {
        delete beverage;
    }
};
\end{lstlisting}

\pause
The decorator wraps another Beverage object.

Every decorator stores a pointer to another Beverage.

\end{frame}

%%%%%%%%%%%%%%%%%%%%%%%%%%%%%%%%%%%%%%%%%%%%%%%%%%%%%
\begin{frame}[fragile]{Concrete Decorator: Milk}

\begin{lstlisting}
class Milk : public CondimentDecorator {
public:
    Milk(Beverage* bev) : CondimentDecorator(bev) {}

    string getDescription() const override {
        return beverage->getDescription() + ", Milk";
    }

    double cost() const override {
        return beverage->cost() + 5.0;
    }
};
\end{lstlisting}

\pause

\end{frame}

%%%%%%%%%%%%%%%%%%%%%%%%%%%%%%%%%%%%%%%%%%%%%%%%%%%%%

\begin{frame}[fragile]{Concrete Decorator: Mocha}

\begin{lstlisting}
class Mocha : public CondimentDecorator {
public:
    Mocha(Beverage* bev) : CondimentDecorator(bev) {}
    string getDescription() const override { return beverage->getDescription() + ", Mocha";
    }

    double cost() const override {
        return beverage->cost() + 8.0;
    }
};
\end{lstlisting}

\pause

Every new decorator extends the previous Beverage.

No existing class needs modification.

\end{frame}

%%%%%%%%%%%%%%%%%%%%%%%%%%%%%%%%%%%%%%%%%%%%%%%%%%%%%

\begin{frame}[fragile]{Client Code}

\begin{lstlisting}
int main() {
    Beverage* myDrink = new DarkRoast();

    myDrink = new Milk(myDrink);
    myDrink = new Mocha(myDrink);

    cout << myDrink->getDescription() << endl;
    cout << myDrink->cost();
    delete myDrink;
}
\end{lstlisting}

\end{frame}

%%%%%%%%%%%%%%%%%%%%%%%%%%%%%%%%%%%%%%%%%%%%%%%%%%%%%

\begin{frame}{How Decoration Works}

The object is wrapped layer by layer.

\bigskip

\[
DarkRoast
\]

\[
\Uparrow
\]

\[
Milk(DarkRoast)
\]

\[
\Uparrow
\]

\[
Mocha(Milk(DarkRoast))
\]

\pause

Every wrapper is still a Beverage.

\end{frame}

%%%%%%%%%%%%%%%%%%%%%%%%%%%%%%%%%%%%%%%%%%%%%%%%%%%%%

\begin{frame}{Execution Flow}

Calling

\[
myDrink->cost()
\]

starts from the outermost decorator.

\bigskip

Execution order:

\[
Mocha
\rightarrow
Milk
\rightarrow
DarkRoast
\]

Then the values are returned back.

\[
8
+5
+30
=
43
\]

\end{frame}

%%%%%%%%%%%%%%%%%%%%%%%%%%%%%%%%%%%%%%%%%%%%%%%%%%%%%

\begin{frame}{Program Output}

\Large

Description

\[
\boxed{
Dark\ Roast\ Coffee,\ Milk,\ Mocha
}
\]

Total Cost

\[
\boxed{43k\ VND}
\]

\end{frame}

%%%%%%%%%%%%%%%%%%%%%%%%%%%%%%%%%%%%%%%%%%%%%%%%%%%%%

\section{Advantages and Disadvantages}

\begin{frame}{Advantages}

\begin{block}{Runtime Flexibility}

Features can be added dynamically while the program is running.

\end{block}

\pause

\begin{block}{Unlimited Combinations}

Supports any topping combination without creating new subclasses.

\end{block}

\pause

\begin{block}{Reusable Components}

Each decorator can be reused independently.

\end{block}

\pause

\begin{block}{Follows SOLID}

\begin{itemize}

\item Single Responsibility Principle

\item Open/Closed Principle

\end{itemize}

\end{block}

\end{frame}

%%%%%%%%%%%%%%%%%%%%%%%%%%%%%%%%%%%%%%%%%%%%%%%%%%%%%

\begin{frame}{Disadvantages}

\begin{block}{Many Small Classes}

Each feature requires its own decorator class.

\end{block}

\pause

\begin{block}{Harder Debugging}

Many wrapper layers produce long call stacks.

\end{block}

\pause

\begin{block}{Complex Initialization}

Nested construction may become difficult.

Example:

\[
new\ Mocha(
new\ Milk(
new\ Mocha(
new\ DarkRoast())))
\]

\end{block}

\end{frame}

\section{Other Patterns}

\begin{frame}{Strategy Pattern}
    \begin{block}{Core Idea}
        Encapsulate different algorithms behind a common interface
        and allow them to be switched at runtime.
    \end{block}

    \begin{itemize}
        \item One task $\rightarrow$ Multiple possible algorithms
        \item Each algorithm is implemented in a separate class
        \item The algorithm can be changed dynamically at runtime
        \item Avoids large \texttt{if-else} or \texttt{switch-case} blocks
    \end{itemize}

    \textbf{Example:} A graph visualizer supporting multiple pathfinding algorithms.
\end{frame}

%=========================================================
\begin{frame}[fragile]{The Problem: Without Strategy}

\begin{lstlisting}
class GraphVisualizer {
public:
    void RunPathfinding(string algoType) {
        if (algoType == "Dijkstra") {
        }
        else if (algoType == "BellmanFord") {
        }
        else if (algoType == "BFS") {
        }
    }
};
\end{lstlisting}

\begin{alertblock}{Problems}
    \begin{itemize}
        \item The class becomes very large and difficult to maintain
        \item Adding a new algorithm requires modifying existing code
        \item High risk of introducing bugs
        \item Violates the Open-Closed Principle
    \end{itemize}
\end{alertblock}

\end{frame}

%=========================================================
\begin{frame}[fragile]{The Solution: Strategy Pattern}

\begin{lstlisting}
class IPathfindingStrategy {
public:
    virtual void ExecuteAlgorithm() = 0;
    virtual ~IPathfindingStrategy() {}
};
\end{lstlisting}

\begin{block}{Common Interface}
    Every pathfinding algorithm follows the same contract.
\end{block}

\begin{itemize}
    \item \texttt{DijkstraStrategy}
    \item \texttt{BellmanFordStrategy}
    \item \texttt{BFSStrategy}
    \item \texttt{...}
\end{itemize}

Each algorithm is encapsulated in its own class.
The main application does not need to know the implementation details.
\end{frame}

%=========================================================
\begin{frame}[fragile]{Concrete Strategies}

\begin{lstlisting}
class DijkstraStrategy : public IPathfindingStrategy {
public:
    void ExecuteAlgorithm() override { // Dijkstra's algorithm
    }
};

class BellmanFordStrategy : public IPathfindingStrategy {
public:
    void ExecuteAlgorithm() override { // Bellman-Ford algorithm
    }
};
\end{lstlisting}

\begin{block}{Key Idea}
    Different algorithms are interchangeable because
    they implement the same interface.
\end{block}

\end{frame}

%=========================================================
\begin{frame}[fragile]{Context: Graph Visualizer}

\begin{lstlisting}
class GraphVisualizer {
private:
    IPathfindingStrategy* strategy;
public:
    void SetStrategy(IPathfindingStrategy* s) {
        strategy = s;
    }
    void Run() {
        if (strategy)
            strategy->ExecuteAlgorithm();
    }
};
\end{lstlisting}

\begin{block}{Context}
    The \texttt{GraphVisualizer} does not care which algorithm is used.
    It simply executes the selected strategy.
\end{block}

\end{frame}

%=========================================================
\begin{frame}[fragile]{Switching Strategies at Runtime}

\begin{lstlisting}
GraphVisualizer app;

// User selects Dijkstra
app.SetStrategy(new DijkstraStrategy());
app.Run();

// User switches to Bellman-Ford
app.SetStrategy(new BellmanFordStrategy());
app.Run();
\end{lstlisting}

\begin{alertblock}{Result}
    The algorithm changes, but the \texttt{GraphVisualizer} code
    remains unchanged.
\end{alertblock}

\end{frame}

%=========================================================
\begin{frame}{Strategy Pattern: Summary}

\begin{center}
    \Large
    \textbf{One task, multiple interchangeable algorithms}
\end{center}

\begin{itemize}
    \item Encapsulate each algorithm separately
    \item Use a common interface
    \item Switch algorithms at runtime
    \item Reduce complex conditional logic
    \item Easy to extend with new strategies
\end{itemize}

\begin{block}{Main Benefit}
    Add a new algorithm without modifying the existing context.
\end{block}

\end{frame}


%=========================================================
\section{Facade Pattern}

\begin{frame}{Facade Pattern}

    \begin{block}{Core Idea}
        Provide a simple, unified interface to a complex subsystem.
    \end{block}

    \begin{center}
        \Large
        \textbf{Complex System}
        $\quad\longrightarrow\quad$
        \textbf{Simple Interface}
    \end{center}

    \begin{itemize}
        \item Hide unnecessary implementation details
        \item Reduce the complexity seen by the client
        \item Provide a single entry point to a subsystem
    \end{itemize}

    \textbf{Example:} Starting a computer by pressing the power button.
\end{frame}

%=========================================================
\begin{frame}[fragile]{The Problem: A Complex Subsystem}

\begin{lstlisting}
class CPU {
public:
    void Freeze() {}
    void Jump(long position) {}
    void Execute() {}
};
class Memory {
public:
    void Load(long position, char* data) {}
};
class HardDrive {
public:
    char* Read(long lba, int size) {
        return "OS data";
    }
};
\end{lstlisting}
\end{frame}

%=========================================================
\begin{frame}{The Problem: A Complex Subsystem (cont.)}
\begin{block}{Problem}
    Starting a computer requires coordinating multiple components:
    CPU, RAM, and Hard Drive.
\end{block}

\end{frame}

%=========================================================
\begin{frame}[fragile]{The Solution: Computer Facade}

\begin{lstlisting}
class ComputerFacade {
private:
    CPU cpu;
    Memory ram;
    HardDrive hdd;

public:
    void StartComputer() {
        cpu.Freeze();
        ram.Load(0x00, hdd.Read(0, 1024));
        cpu.Jump(0x00);
        cpu.Execute();
    }
};
\end{lstlisting}
\end{frame}

\begin{frame}{Facade Pattern}
\begin{block}{Facade}
    \texttt{ComputerFacade} hides the complex startup process
    behind a single method: \texttt{StartComputer()}.
\end{block}
\end{frame}


%=========================================================
\begin{frame}[fragile]{Client: Simple Interface}

\begin{lstlisting}
int main() {
    ComputerFacade myPC;

    // Press the power button
    myPC.StartComputer();

    return 0;
}
\end{lstlisting}

\begin{center}
    \Large
    \textbf{One method call. That's it.}
\end{center}

\begin{itemize}
    \item The client does not interact directly with CPU, RAM, or HDD
    \item Internal complexity is hidden
    \item The client only depends on the Facade
\end{itemize}

\end{frame}

%=========================================================
\begin{frame}{Facade Pattern: Real-World Analogy}

\begin{center}
    \Large
    \textbf{Press the Power Button}
\end{center}

\begin{center}
    \texttt{StartComputer()}
\end{center}

\begin{itemize}
    \item CPU initialization
    \item Loading data from the Hard Drive
    \item Loading data into Memory
    \item Starting CPU execution
\end{itemize}

\begin{block}{The User}
    Only needs to press one button.
    They do not need to understand how the internal components work together.
\end{block}

\end{frame}

%=========================================================
\section{Relationships with Other Design Patterns}

\begin{frame}{Relationships with Other Design Patterns}

\begin{center}
    \Large
    \textbf{Decorator, Facade, and Strategy solve different problems}
\end{center}

\vspace{0.5cm}

\begin{columns}[T]

% Decorator
\column{0.32\textwidth}
\begin{block}{Decorator}
    \centering
    \textbf{Add Behavior}
    
    \vspace{0.3cm}
    
    Extends an object's functionality
    without modifying its original class.
\end{block}

% Facade
\column{0.32\textwidth}
\begin{block}{Facade}
    \centering
    \textbf{Hide Complexity}
    
    \vspace{0.3cm}
    
    Provides a simple interface
    to a complex subsystem.
\end{block}

% Strategy
\column{0.32\textwidth}
\begin{block}{Strategy}
    \centering
    \textbf{Change Behavior}
    
    \vspace{0.3cm}
    
    Allows different algorithms
    to be switched at runtime.
\end{block}

\end{columns}

\vspace{0.6cm}

\begin{center}
    \Large
    \textbf{Decorator} $\rightarrow$ \textit{What can I add?}
    
    \textbf{Facade} $\rightarrow$ \textit{How can I simplify access?}
    
    \textbf{Strategy} $\rightarrow$ \textit{How can I change the algorithm?}
\end{center}

\end{frame}


%=========================================================
\begin{frame}{Decorator, Facade, and Strategy}

\begin{center}
\begin{tabular}{|p{2.8cm}|p{3.5cm}|p{3.5cm}|}
\hline
\textbf{Pattern} & \textbf{Main Purpose} & \textbf{Key Question} \\
\hline

\textbf{Decorator}
&
Add or extend behavior
&
\textit{How can I add functionality?}
\\
\hline

\textbf{Facade}
&
Simplify access to a complex system
&
\textit{How can I make the system easier to use?}
\\
\hline

\textbf{Strategy}
&
Switch between different algorithms
&
\textit{How can I change the way a task is performed?}
\\
\hline

\end{tabular}
\end{center}

\vspace{0.6cm}

\begin{block}{The Big Picture}
    All three patterns help reduce coupling and improve flexibility,
    but they focus on different aspects of software design.
\end{block}

\end{frame}


%=========================================================
\begin{frame}{How They Can Work Together}

\begin{center}
    \Large
    \textbf{Example: A Pathfinding System}
\end{center}

\begin{block}{Solution}
    The client calls one simple method through the Facade.
    The Strategy selects the pathfinding algorithm,
    while Decorators add extra functionality.
\end{block}

\end{frame}


%=========================================================
\begin{frame}{Summary}

\begin{center}

\vspace{0.5cm}

\Large
\textbf{Decorator}
\\
\large
Adds new behavior to an object.

\vspace{0.5cm}

\textbf{Facade}
\\
\large
Hides system complexity behind a simple interface.

\vspace{0.5cm}

\textbf{Strategy}
\\
\large
Allows different algorithms to be exchanged easily.

\vspace{0.8cm}

\begin{block}{Key Takeaway}
\centering
\textbf{
Decorator adds behavior,
Facade simplifies access,
Strategy changes behavior.
}
\end{block}

\end{center}

\end{frame}

%%%%%%%%%%%%%%%%%%%%%%%%%%%%%%%%%%%%%%%%%%%%%%%%%%%%%
\section{Real-world Applications}

%%%%%%%%%%%%%%%%%%%%%%%%%%%%%%%%%%%%%%%%%%%%%%%%%%%%%
% DOMAIN 1: JAVA I/O STREAMS
%%%%%%%%%%%%%%%%%%%%%%%%%%%%%%%%%%%%%%%%%%%%%%%%%%%%%

\begin{frame}{9. Real-World Applications: Java I/O Streams}

\begin{block}{Conceptual Background}
In the \texttt{java.io} package, reading data can originate from multiple sources (files, network sockets, byte arrays) and require various processing capabilities (buffering, decompression, data parsing). Creating subclass combinations like \texttt{BufferedGzipFileInputStream} would cause severe class explosion.
\end{block}

\pause

\begin{block}{GoF Design Pattern Mapping}
\begin{itemize}
    \item \textbf{Component}: \texttt{InputStream} (Abstract base interface defining \texttt{read()})
    \item \textbf{Concrete Component}: \texttt{FileInputStream} (Reads raw bytes directly from a file)
    \item \textbf{Base Decorator}: \texttt{FilterInputStream} (Holds an internal \texttt{InputStream} reference)
    \item \textbf{Concrete Decorators}: \texttt{BufferedInputStream}, \texttt{GZIPInputStream} (Add buffering and decompression)
\end{itemize}
\end{block}

\end{frame}

%%%%%%%%%%%%%%%%%%%%%%%%%%%%%%%%%%%%%%%%%%%%%%%%%%%%%

\begin{frame}[fragile]{Java I/O Streams: Execution Flow}

\begin{lstlisting}[language=Java]
// Base file reader wrapped with buffering and GZIP decompression
InputStream fileStream = new FileInputStream("data.txt.gz");
InputStream bufferedStream = new BufferedInputStream(fileStream);
InputStream gzipStream = new GZIPInputStream(bufferedStream);

int data = gzipStream.read(); // Execution propagates down through wrappers
\end{lstlisting}

\pause

\begin{block}{How Decoration Works at Runtime}
\begin{enumerate}
    \item The client calls \texttt{read()} on the outermost decorator (\texttt{GZIPInputStream}).
    \item \texttt{GZIPInputStream} requests raw compressed chunks from \texttt{BufferedInputStream}.
    \item \texttt{BufferedInputStream} fetches byte blocks from \texttt{FileInputStream}.
    \item Data flows back up through the stack, being buffered and decompressed dynamically!
\end{enumerate}
\end{block}

\end{frame}

%%%%%%%%%%%%%%%%%%%%%%%%%%%%%%%%%%%%%%%%%%%%%%%%%%%%%
% DOMAIN 2: BACKEND MIDDLEWARE
%%%%%%%%%%%%%%%%%%%%%%%%%%%%%%%%%%%%%%%%%%%%%%%%%%%%%

\begin{frame}{Backend Web Middleware}

\begin{block}{Architectural Problem}
Modern backend frameworks (Express.js, ASP.NET Core, Spring Filters) must process incoming HTTP requests with cross-cutting concerns like authentication, CORS headers, logging, and error handling before reaching the final controller.
\end{block}

\pause

\begin{block}{The Middleware Wrapper Concept}
Instead of polluting every controller endpoint with repetitive security checks, middleware functions act as an onion-skin wrapper stack around the business logic controller.
\end{block}

\begin{center}
\small
\texttt{Request} $\rightarrow$ \boxed{\text{Logger}} $\rightarrow$ \boxed{\text{AuthCheck}} $\rightarrow$ \boxed{\text{CorsHeader}} $\rightarrow$ \texttt{Controller}
\end{center}

\end{frame}

%%%%%%%%%%%%%%%%%%%%%%%%%%%%%%%%%%%%%%%%%%%%%%%%%%%%%

\begin{frame}[fragile]{Backend Middleware: Code \& Delegation}

\begin{lstlisting}[language=C++]
// Conceptual Express.js / ASP.NET Middleware Stack
void AuthMiddleware(Request req, Response res, NextFunction next) {
    if (!req.hasValidToken()) {
        res.sendError(401, "Unauthorized");
        return; // Intercept execution early
    }
    next(); // Pass control to the next decorated layer in the pipeline
}
\end{lstlisting}

\pause

\begin{block}{Why Decorator Pattern Fits Here}
\begin{itemize}
    \item \textbf{Dynamic Stacking}: Middleware layers can be added or removed per route without modifying the target route controller.
    \item \textbf{Pre/Post Processing}: A middleware layer can execute code both \textit{before} calling \texttt{next()} and \textit{after} \texttt{next()} returns (e.g., measuring response latency).
\end{itemize}
\end{block}

\end{frame}

%%%%%%%%%%%%%%%%%%%%%%%%%%%%%%%%%%%%%%%%%%%%%%%%%%%%%
% DOMAIN 3: PYTHON DECORATORS
%%%%%%%%%%%%%%%%%%%%%%%%%%%%%%%%%%%%%%%%%%%%%%%%%%%%%

\begin{frame}{Python Decorators}

\begin{block}{Native Language Support}
Python integrates the Decorator pattern directly into its language syntax via the \texttt{@decorator\_name} annotation. In Python, functions are first-class objects, allowing higher-order functions to wrap target callables dynamically.
\end{block}

\pause

\begin{block}{Syntactic Sugar Mechanism}
Writing an annotation in Python:
\begin{center}
\texttt{@login\_required}\\
\texttt{def view\_dashboard(): ...}
\end{center}
is exact syntactic sugar for wrapping the function object at definition time:
\begin{center}
\texttt{view\_dashboard = login\_required(view\_dashboard)}
\end{center}
\end{block}

\end{frame}

%%%%%%%%%%%%%%%%%%%%%%%%%%%%%%%%%%%%%%%%%%%%%%%%%%%%%

\begin{frame}[fragile]{Python Decorators: Practical Implementation}

\begin{lstlisting}[language=Python]
def login_required(func):
    def wrapper(*args, **kwargs):
        if not current_user.is_authenticated:
            raise PermissionError("User must log in to view page")
        return func(*args, **kwargs) # Delegate to original function
    return wrapper

@login_required
def view_dashboard():
    return "Secret User Dashboard Data"
\end{lstlisting}

\pause

\begin{block}{Key Takeaway}
The \texttt{login\_required} wrapper adds authentication checks dynamically without modifying a single line of code inside \texttt{view\_dashboard()}.
\end{block}

\end{frame}

%%%%%%%%%%%%%%%%%%%%%%%%%%%%%%%%%%%%%%%%%%%%%%%%%%%%%
% DOMAIN 4: REACT HIGHER-ORDER COMPONENTS (HOC)
%%%%%%%%%%%%%%%%%%%%%%%%%%%%%%%%%%%%%%%%%%%%%%%%%%%%%

\begin{frame}{React Higher-Order Components (HOC)}

\begin{block}{UI Component Decoration}
In modern web frontend development with React, user interface components should focus on rendering layout. Cross-cutting UI concerns (such as authorization checks, global theme injection, or data fetching) are attached using Higher-Order Components.
\end{block}

\pause

\begin{block}{GoF Structural Mapping in React}
\begin{itemize}
    \item \textbf{Component}: Base React Component (\texttt{UserProfilePage})
    \item \textbf{Decorator HOC}: Wrapper function (\texttt{withAuth}) returning an enhanced Component
    \item \textbf{Concrete Result}: \texttt{withAuth(UserProfilePage)}
\end{itemize}
\end{block}

\end{frame}

%%%%%%%%%%%%%%%%%%%%%%%%%%%%%%%%%%%%%%%%%%%%%%%%%%%%%

\begin{frame}[fragile]{React HOC: Code \& Composition}

\begin{lstlisting}[language=Java]
// React Higher-Order Component Wrapper Example
const withAuth = (WrappedComponent) => {
    return (props) => {
        if (!props.isAuthenticated) {
            return <LoginPage />;
        }
        return <WrappedComponent {...props} />;
    };
};

// Usage: Create an authenticated view without touching UserProfilePage source code
const ProtectedProfile = withAuth(UserProfilePage);
\end{lstlisting}

\pause

\begin{block}{Benefit}
Multiple HOC decorators can be composed: \texttt{withTheme(withAuth(UserProfilePage))}, giving UI components modular, reusable capabilities.
\end{block}

\end{frame}

%%%%%%%%%%%%%%%%%%%%%%%%%%%%%%%%%%%%%%%%%%%%%%%%%%%%%
% DOMAIN 5: FLUTTER WIDGET COMPOSITION
%%%%%%%%%%%%%%%%%%%%%%%%%%%%%%%%%%%%%%%%%%%%%%%%%%%%%

\begin{frame}{Flutter Widget Trees}

\begin{block}{Composition Over Inheritance in Mobile UI}
Google's Flutter framework completely abandons traditional subclassing for UI elements. Instead of creating a \texttt{PaddedBorderedText} subclass, Flutter adopts the philosophy: \textbf{"Everything is a Widget"}.
\end{block}

\pause

\begin{block}{Visual Decorator Tree}
Base content is decorated visually by nesting wrapper widgets layer by layer:
\begin{center}
\texttt{Container} (Adds background color \& border)\\
$\uparrow$\\
\texttt{Padding} (Adds outer spacing)\\
$\uparrow$\\
\texttt{Text("Hello CS202")} (Base visual component)
\end{center}
\end{block}

\end{frame}

%%%%%%%%%%%%%%%%%%%%%%%%%%%%%%%%%%%%%%%%%%%%%%%%%%%%%

\begin{frame}[fragile]{Flutter Widgets: Code \& Render Pipeline}

\begin{lstlisting}[language=Java]
// Flutter UI composition via nested decorator widgets
Widget build(BuildContext context) {
    return Container(
        color: Colors.blue,
        child: Padding(
            padding: EdgeInsets.all(16.0),
            child: Text("Hello CS202 Seminar"),
        ),
    );
}
\end{lstlisting}

\pause

\begin{block}{Render Execution Flow}
During rendering, Flutter passes layout constraints down through the \texttt{Container} and \texttt{Padding} wrappers to the inner \texttt{Text} child. The computed pixel dimensions return back up the tree-exactly mirroring how decorator cost calculations bubble up in object wrapping!
\end{block}

\end{frame}}

%%%%%%%%%%%%%%%%%%%%%%%%%%%%%%%%%%%%%%%%%%%%%%%%%%%%%
\section{Mini Quiz}

\begin{frame}{Question 1}

Decorator belongs to which Design Pattern category?

\bigskip

A. Creational

B. Structural

C. Behavioral

D. Architectural

\pause

\begin{alertblock}{Answer}

B. Structural

\end{alertblock}

\end{frame}

%%%%%%%%%%%%%%%%%%%%%%%%%%%%%%%%%%%%%%%%%%%%%%%%%%%%%

\begin{frame}{Question 2}

Which design principle does Decorator emphasize?

\bigskip

A. Inheritance over Composition

B. Composition over Inheritance

C. Tight Coupling

D. Procedural Programming

\pause

\begin{alertblock}{Answer}

B. Composition over Inheritance

\end{alertblock}

\end{frame}

%%%%%%%%%%%%%%%%%%%%%%%%%%%%%%%%%%%%%%%%%%%%%%%%%%%%%

\begin{frame}{Question 3}

What is the main purpose of Decorator?

\bigskip

A. Clone an object

B. Restrict object creation

C. Add new behavior at Runtime

D. Simplify a subsystem

\pause

\begin{alertblock}{Answer}

C. Add new behavior at Runtime

\end{alertblock}

\end{frame}

%%%%%%%%%%%%%%%%%%%%%%%%%%%%%%%%%%%%%%%%%%%%%%%%%%%%%

\begin{frame}{Question 4}

Which is a disadvantage of Decorator?

\bigskip

A. Violates OCP

B. Causes Class Explosion

C. Long call stack and many small classes

D. Cannot work in OOP languages

\pause

\begin{alertblock}{Answer}

C. Long call stack and many small classes

\end{alertblock}

\end{frame}

%%%%%%%%%%%%%%%%%%%%%%%%%%%%%%%%%%%%%%%%%%%%%%%%%%%%%

\begin{frame}{Question 5}

How do we implement Double Milk?

\bigskip

A. Create DoubleMilkDecorator

B. Add count = 2

C. Wrap Milk twice

D. Copy the code twice

\pause

\begin{alertblock}{Answer}

C. Wrap Milk twice

\end{alertblock}

\end{frame}

%%%%%%%%%%%%%%%%%%%%%%%%%%%%%%%%%%%%%%%%%%%%%%%%%%%%%

\begin{frame}

\centering

\Huge

Thank you for listening!

\end{frame}

%%%%%%%%%%%%%%%%%%%%%%%%%%%%%%%%%%%%%%%%%%%%%%%%%%%%%

\end{document}